\documentclass[11pt]{article}
\usepackage[margin=1in]{geometry}
\usepackage{amsmath,amssymb,booktabs,array,longtable}
\usepackage[hidelinks]{hyperref}

\title{Search, Selection, Low-Rank Wage Interfaces, and Edge Diagnostics}
\author{}
\date{\today}

\newcommand{\E}{\mathbb{E}}
\newcommand{\1}{\mathbf{1}}
\newcommand{\Var}{\operatorname{Var}}
\newcommand{\Cov}{\operatorname{Cov}}

\begin{document}
\maketitle

% ============================================================
% Merged from ss_to_bs_memo.tex
% ============================================================


\section{From Shimer--Smith to Borovi\v{c}kov\'a--Shimer}

\paragraph{Source convention.}

Page citations to Shimer--Smith refer to the printed Econometrica pages in
\emph{Assortative Matching and Search}. Page citations to Borovi\v{c}kov\'a--Shimer
refer to the paper page numbers printed in the October 23, 2025 draft,
\emph{Assortative Matching and Wages: The Role of Selection}. When I reuse a
result from the prior project notes, I cite the prior result PDF
\texttt{7\_15\_Meeting (1).pdf} and the source \TeX{} files
\texttt{shimer\_smith\_writeup.tex} or
\texttt{model\_acceptance\_stochasticity\_writeup.tex}. Displays or claims not
directly sourced are labeled as derived comparisons or derived identities.

\subsection{Notation}

\begin{longtable}{>{\raggedright\arraybackslash}p{0.24\linewidth}
                >{\raggedright\arraybackslash}p{0.68\linewidth}}
\toprule
Symbol & Meaning \\
\midrule
\endfirsthead
\toprule
Symbol & Meaning \\
\midrule
\endhead
$x,y$ & Types. In Shimer--Smith, both sides are symmetric agent types; in BS, $x$ is a worker type and $y$ is a firm/job type. Source: Shimer--Smith, p. 346; BS, pp. 4--5. \\
$u(y),u_x$ & Unmatched density or measure. Shimer--Smith use $u(y)$ for unmatched type-$y$ agents; BS use $u_x$ for unemployed type-$x$ workers. Source: Shimer--Smith, p. 347; BS, p. 5. \\
$v_y$ & Measure of vacant type-$y$ firms/jobs in BS. Source: BS, p. 5. \\
$\rho$ & Meeting/contact-rate parameter. Source: Shimer--Smith, p. 347; BS, p. 5. \\
$r,\delta$ & Discount rate and exogenous match-separation rate. Source: Shimer--Smith, pp. 347--348; BS, pp. 4--5. \\
$f(x,y),f_{x,y}$ & Type-pair production surface. Shimer--Smith output is deterministic $f(x,y)$; BS output is $zf_{x,y}$. Source: Shimer--Smith, pp. 346--347; BS, p. 5. \\
$z,G,g,z_0,\theta$ & BS match-specific productivity shock, its distribution and density, and the Pareto scale and shape parameters used in Proposition 2. For the Pareto mean formulas with $(\theta-1)$ in the denominator, the relevant assumption is $\theta>1$; a Pareto distribution itself only requires a positive shape parameter. Source: BS, pp. 5, 11. \\
$W(x),W(x\mid y)$ & Shimer--Smith values of being unmatched and matched with type $y$. Source: Shimer--Smith, p. 348. \\
$S(x\mid y)$ & Shimer--Smith personal capital surplus, $W(x\mid y)-W(x)$. Source: Shimer--Smith, p. 348. \\
$\mathcal M(x)$ & Shimer--Smith mutually agreeable matching set for type $x$. Source: Shimer--Smith, p. 347. \\
$w(x)$ & Shimer--Smith reservation wage / flow value of being unmatched, $w(x)=rW(x)$. Source: Shimer--Smith, p. 349. \\
$\pi(x\mid y)$ & Shimer--Smith flow payoff to type $x$ while matched with type $y$. Source: Shimer--Smith, pp. 347--348. \\
$\alpha_x,\psi_y,c_{xy}$ & Derived additive worker-side term, additive other-side term, and interaction term used to rewrite the Shimer--Smith payoff as $\pi(x\mid y)=\alpha_x+\psi_y+c_{xy}$. Source for derivation: prior result PDF, Sec. 6.3; \texttt{shimer\_smith\_writeup.tex}, Sec. 6.3. \\
$V_x^u,V_y^v$ & BS values of unemployment for worker type $x$ and vacancy for firm type $y$. Source: BS, pp. 5--6. \\
$V^e_{x,y}(z,w),V^f_{y,x}(z,w)$ & BS employed-worker and matched-firm values for match quality $z$ and wage $w$. Source: BS, pp. 5--6. \\
$\bar w_x,\bar\pi_y$ & BS worker reservation wage and firm reservation profit, $\bar w_x=rV_x^u$ and $\bar\pi_y=rV_y^v$. Source: BS, p. 6. \\
$V^s_{x,y}(z)$ & BS total match surplus. Source: BS equation (5), p. 6. \\
$\bar z_{x,y}$ & BS reservation productivity threshold, $(\bar w_x+\bar\pi_y)/f_{x,y}$. Source: BS equation (6), p. 6. \\
$\gamma$ & BS worker bargaining weight in Nash bargaining. Source: BS, p. 5. \\
$W_{x,y}(z)$ & BS equilibrium wage in an accepted match. Source: BS equation (9), p. 7. \\
$ITT_x(y),ATT_x(y)$ & BS intent-to-treat and average treatment on the treated from creating a meeting with a type-$y$ firm for a type-$x$ worker. Source: BS equations (17)--(19), pp. 9--10. \\
$\omega_{i,j},\alpha_x^*,\psi_y^*,\epsilon_{i,j}$ & BS observed accepted-match wage and Proposition 2 worker component, firm component, and residual. Source: BS Proposition 2, pp. 11--12. \\
\bottomrule
\end{longtable}

\subsection{The Clean Mapping}

The old notes had the right informal comparison: Shimer--Smith is random search
with deterministic type-pair surplus, while Borovi\v{c}kov\'a--Shimer is random
search with stochastic match-specific surplus. Source: prior result PDF, Secs.
1 and 8.4; \texttt{model\_acceptance\_stochasticity\_writeup.tex}, Secs. 3--4.
The mapping is:

\[
\begin{array}{>{\raggedright\arraybackslash}p{0.27\linewidth}
                >{\raggedright\arraybackslash}p{0.31\linewidth}
                >{\raggedright\arraybackslash}p{0.31\linewidth}}
\toprule
Object & Shimer--Smith & Borovi\v{c}kov\'a--Shimer \\
\midrule
Side structure & Symmetric agents of types $x,y$ & Workers $x$ and firms/jobs $y$ \\
Meetings & Unmatched type $x$ meets unmatched type $y$ at rate
$\rho u(y)\,dy$ & Unemployed type $x$ worker contacts vacant type $y$ firm at
rate $\rho v_y$ \\
Match output & Deterministic $f(x,y)$ & Stochastic $z f_{x,y}$, with $z\sim G$ drawn at the meeting \\
Worker outside flow value & $w(x)=rW(x)$ & $\bar w_x=rV_x^u$ \\
Other side outside flow value & $w(y)=rW(y)$ & $\bar\pi_y=rV_y^v$ \\
Accepted-match payoff/wage & $\pi(x\mid y)$, the worker's flow payoff & $W_{x,y}(z)$, the equilibrium wage \\
Surplus split & Symmetric Nash, $S(x\mid y)=S(y\mid x)$ & Generalized Nash, worker share $\gamma$ \\
Acceptance & $f(x,y)\geq w(x)+w(y)$ & $z f_{x,y}\geq \bar w_x+\bar\pi_y$ \\
\bottomrule
\end{array}
\]

The following correspondence is a derived notational map across the two models:
\[
f(x,y)
\quad\longleftrightarrow\quad
z f_{x,y},
\qquad
w(x)
\quad\longleftrightarrow\quad
\bar w_x,
\qquad
w(y)
\quad\longleftrightarrow\quad
\bar\pi_y.
\]
The way to explain the table is in three steps.

\begin{enumerate}
\item \textbf{Both papers have random meetings.} In Shimer--Smith, the randomness is
that an unmatched agent does not choose exactly whom to meet; meetings arrive through
the search process. In BS, unemployed workers and vacancies also meet randomly. Source:
Shimer--Smith, p. 347; BS, p. 5.

\item \textbf{The production object differs.} In Shimer--Smith, once the two types
$(x,y)$ are known, the match output is the fixed number $f(x,y)$. In BS, once the
worker type $x$ and firm type $y$ are known, the meeting still draws $z$, so output is
the random value $zf_{x,y}$. Source: Shimer--Smith, pp. 346--347; BS, p. 5.

\item \textbf{The acceptance object therefore differs.} In Shimer--Smith, a type pair is
either acceptable or not acceptable in equilibrium. In BS, the same type pair can sometimes
match and sometimes fail to match, depending on the realized $z$. This is the selection
channel.
\end{enumerate}

Derived comparison: here is the explicit collapse to the Shimer--Smith payoff
formula. Start from the BS accepted-match wage equation:
\[
W_{x,y}(z)
=
\bar w_x+\gamma\left(zf_{x,y}-\bar w_x-\bar\pi_y\right),
\qquad z\geq \bar z_{x,y}.
\]
Source: BS equation (9), p. 7. If bargaining is symmetric, $\gamma=1/2$, then
\[
W_{x,y}(z)
=
\bar w_x+\frac12 zf_{x,y}
-\frac12\bar w_x-\frac12\bar\pi_y
=
\frac12 zf_{x,y}
+\frac12\bar w_x
-\frac12\bar\pi_y.
\]
Now impose the degenerate-productivity case $z=z^\ast$, so every meeting of the
same type pair has the same realized output. Define the deterministic type-pair
production term
\[
\tilde f(x,y)\equiv z^\ast f_{x,y}.
\]
Then the BS wage becomes
\[
W_{x,y}
=
\frac12 \tilde f(x,y)
+\frac12\bar w_x
-\frac12\bar\pi_y.
\tag{BS-degenerate}
\]
Compare this with the Shimer--Smith payoff formula derived in the old notes:
\[
\pi(x\mid y)
=
\frac12 f(x,y)
+\frac12 w(x)
-\frac12 w(y).
\tag{SS-payoff}
\]
Source: prior result PDF, Sec. 6.3;
\texttt{shimer\_smith\_writeup.tex}, Sec. 6.3. The correspondence is therefore
\[
\pi(x\mid y)\leftrightarrow W_{x,y},
\qquad
f(x,y)\leftrightarrow \tilde f(x,y)=z^\ast f_{x,y},
\qquad
w(x)\leftrightarrow \bar w_x,
\qquad
w(y)\leftrightarrow \bar\pi_y.
\]
So, under symmetric bargaining and degenerate $z$, BS has the same algebraic
payoff shape as Shimer--Smith, except that BS labels the two sides as worker and
firm rather than as two symmetric agents.

The reason this shortcut fails with nondegenerate $z$ is also explicit. Holding
the type pair $(x,y)$ fixed, BS no longer has one number $\tilde f(x,y)$. It has
a random realized production value
\[
zf_{x,y}.
\]
Thus the same type pair generates the random wage
\[
W_{x,y}(z)
=
\frac12 zf_{x,y}
+\frac12\bar w_x
-\frac12\bar\pi_y
\quad\text{when }\gamma=1/2,
\]
but only if the draw clears the threshold
\[
z\geq \bar z_{x,y}
=
\frac{\bar w_x+\bar\pi_y}{f_{x,y}}.
\]
Source: BS equation (6), p. 6. Therefore the type pair $(x,y)$ maps to a
distribution of possible realized payoffs, not to one deterministic payoff. The
econometrician sees wages only for the selected meetings whose $z$ clears the
acceptance threshold.

\subsection{Where the Models Actually Diverge}

The divergence is not the Bellman logic. In both papers, value equations have
current flow payoffs plus expected capital gains/losses from Poisson events. In
Shimer--Smith, an unmatched type $x$ receives a capital gain when a mutually
agreeable type $y$ is met, and a matched agent loses the personal surplus when
the match is destroyed. Source: Shimer--Smith, pp. 347--349; prior result PDF,
Secs. 1.3--1.6; \texttt{shimer\_smith\_writeup.tex}, Sec. 1.

This means the first line of explanation should be: ``do not contrast SS and BS as
deterministic versus stochastic search.'' That contrast is too crude. Search is stochastic
in both models. The useful contrast is deterministic type-pair surplus versus stochastic
type-pair surplus.

The divergence is also not simply that one model is random and the other is
deterministic. Both have random meetings and random separations. The exact
divergence is the timing and role of match quality. In Shimer--Smith, once the
types $(x,y)$ are fixed, match output is fixed at $f(x,y)$. Therefore the
equilibrium acceptance rule is deterministic at the type-pair level:
\[
y\in \mathcal M(x)
\quad\Longleftrightarrow\quad
f(x,y)-w(x)-w(y)\geq 0 .
\]
Source: Shimer--Smith, equation (5), p. 349.

In Borovi\v{c}kov\'a--Shimer, a type $x$ worker and type $y$ firm draw $z$ after
they meet. Holding $(x,y)$ fixed, some meetings are accepted and others are
rejected. The type-pair acceptance probability is
\[
\Pr(\text{accept}\mid x,y)=1-G(\bar z_{x,y}),
\qquad
\bar z_{x,y}=\frac{\bar w_x+\bar\pi_y}{f_{x,y}}.
\]
Source: BS equations (6) and (13), pp. 6--8. That is the selection wedge. It
creates a difference between the value of creating a meeting and the value of an
observed accepted match.

The nuance is important. Suppose a worker of type $x$ meets a firm of type $y$ ten
times in ten independent copies of the BS economy. The first-stage type pair is the
same every time, but the $z$ draw changes. Some meetings produce a high enough $z$
and become observed jobs; other meetings are rejected and leave no wage observation.
In Shimer--Smith, by contrast, the same type pair has the same surplus every time,
so the equilibrium type-pair acceptance decision does not vary across repeated meetings.

\subsection{Bellman to Threshold to ITT}

Start with the Shimer--Smith Bellman logic. For an unmatched type $x$ agent,
meetings with mutually acceptable type $y$ agents arrive at rate $\rho u(y)\,dy$.
The value gain from such a meeting is the personal surplus
\[
S(x\mid y)=W(x\mid y)-W(x).
\]
Thus the unmatched Bellman equation is
\[
rW(x)
=
\rho\int_{\mathcal M(x)}S(x\mid y)u(y)\,dy.
\tag{SS-unmatched}
\]
Source: Shimer--Smith, pp. 347--349; prior result PDF, Sec. 1.3;
\texttt{shimer\_smith\_writeup.tex}, Sec. 1.2. The key point is that the
Bellman equation values search by integrating the surplus gains from meetings
that become matches.

Now use the matched Bellman equation. While matched with type $y$, type $x$
receives flow payoff $\pi(x\mid y)$ and loses the match at rate $\delta$. Hence
\[
rW(x\mid y)
=
\pi(x\mid y)-\delta S(x\mid y).
\tag{SS-matched}
\]
Source: Shimer--Smith, pp. 347--349; prior result PDF, Sec. 1.4;
\texttt{shimer\_smith\_writeup.tex}, Sec. 1.3. Since
$S(x\mid y)=W(x\mid y)-W(x)$ and $w(x)=rW(x)$, this can be rearranged as
\[
\pi(x\mid y)
=
w(x)+(r+\delta)S(x\mid y).
\]
Why does this look like only one side of the market? Because Shimer--Smith is a
symmetric matching model. There are not workers and firms. There are agents of
different types, and either side of the match can be called the $x$ side or the
$y$ side. The other side's Bellman equation is therefore not a different economic
object; it is the same equation with $x$ and $y$ swapped:
\[
rW(y\mid x)
=
\pi(y\mid x)-\delta S(y\mid x).
\tag{SS-matched-y}
\]
Under symmetric Nash bargaining in Shimer--Smith,
\[
S(x\mid y)=S(y\mid x),
\]
so the two personal surplus equations are two copies of the same logic. This is
why the notes can write the Bellman equation for type $x$ and then obtain the
type-$y$ version by symmetry.

Using the two side-specific Bellman equations and feasibility,
\[
\pi(x\mid y)+\pi(y\mid x)=f(x,y),
\]
gives the Shimer--Smith personal surplus formula
\[
S(x\mid y)
=
\frac{f(x,y)-w(x)-w(y)}{2(r+\delta)}.
\tag{SS-surplus}
\]
Source for the derivation: prior result PDF, Secs. 1.5--1.6;
\texttt{shimer\_smith\_writeup.tex}, Sec. 1.4. The match is mutually acceptable
exactly when this surplus is nonnegative:
\[
y\in\mathcal M(x)
\quad\Longleftrightarrow\quad
f(x,y)\geq w(x)+w(y).
\tag{SS-threshold}
\]
Source: Shimer--Smith, equation (5), p. 349.

This is the cross-paper bridge. In Shimer--Smith, a meeting with type $y$ has a
deterministic worker-side value
\[
\max\left\{
\frac{f(x,y)-w(x)-w(y)}{2(r+\delta)},0
\right\}.
\]
There is no extra match-quality draw after the meeting, so conditional on the
type pair $(x,y)$, the meeting either always creates this surplus or always
creates zero surplus. In BS, the same Bellman logic is kept, but the meeting
draws $z$ before the acceptance decision. Therefore the value of creating a
meeting must average over the possible $z$ realizations. That average is exactly
the ITT object.

The mathematical connection between the Shimer--Smith matched Bellmans and the
BS Bellmans is a relabeling plus a change in the output object. In Shimer--Smith,
the two side-specific matched equations can be written as
\[
rW(x\mid y)
=
\pi(x\mid y)+\delta\left(W(x)-W(x\mid y)\right),
\tag{SS-x}
\]
\[
rW(y\mid x)
=
\pi(y\mid x)+\delta\left(W(y)-W(y\mid x)\right).
\tag{SS-y}
\]
The two flow payoffs exhaust deterministic output:
\[
\pi(x\mid y)+\pi(y\mid x)=f(x,y).
\]
BS keeps this Bellman form but labels the two sides as worker and firm. The
worker's matched value is $V^e_{x,y}(z,w)$, the worker's outside value is
$V_x^u$, and the worker's flow payoff is the wage $w$. The firm's matched value
is $V^f_{y,x}(z,w)$, the firm's outside value is $V_y^v$, and the firm's flow
payoff is output net of the wage, $zf_{x,y}-w$. Thus the side-by-side map is
\[
\begin{array}{ccl}
W(x\mid y) &\longleftrightarrow& V^e_{x,y}(z,w),\\
W(x) &\longleftrightarrow& V_x^u,\\
\pi(x\mid y) &\longleftrightarrow& w,\\
W(y\mid x) &\longleftrightarrow& V^f_{y,x}(z,w),\\
W(y) &\longleftrightarrow& V_y^v,\\
\pi(y\mid x) &\longleftrightarrow& zf_{x,y}-w,\\
f(x,y) &\longleftrightarrow& zf_{x,y}.
\end{array}
\]

BS write the worker and firm Bellman equations for a realized match quality
$z$ and arbitrary wage $w$ as
\[
rV^e_{x,y}(z,w)=w+\delta\left(V_x^u-V^e_{x,y}(z,w)\right),
\]
\[
rV^f_{y,x}(z,w)=z f_{x,y}-w+\delta\left(V_y^v-V^f_{y,x}(z,w)\right).
\]
Source: BS equations (2)--(4), pp. 6--7. Define
\[
\bar w_x=rV_x^u,\qquad \bar\pi_y=rV_y^v.
\]
This is the BS analogue of the two Shimer--Smith matched Bellman equations:
current flow payoff plus the capital loss from separation. BS writes two
equations because workers and firms are economically distinct sides with
different flow payoffs. Shimer--Smith needs only one displayed side because the
model is symmetric; the other side is obtained by swapping $x$ and $y$.

Adding the worker and firm values and subtracting outside values gives BS total
match surplus
\[
V^s_{x,y}(z)
=
\frac{\max\{z f_{x,y}-\bar w_x-\bar\pi_y,0\}}{r+\delta}.
\]
Source: BS equation (5), p. 6.

Thus the selection threshold is exactly
\[
z\geq \bar z_{x,y},
\qquad
\bar z_{x,y}
=
\frac{\bar w_x+\bar\pi_y}{f_{x,y}}.
\]
Source: BS equation (6), p. 6. This is the BS analogue of Shimer--Smith's
$f(x,y)\geq w(x)+w(y)$, except that the left-hand side is realized only after
the meeting.

Now the ITT follows from the same object that appears inside the Bellman
equation. In Shimer--Smith, the value of a type-$y$ meeting for type $x$ is the
deterministic personal surplus $S(x\mid y)$ when the pair is acceptable. In BS,
the value of a type-$y$ meeting for a type-$x$ worker before seeing $z$ is the
expected worker share of the realized total surplus:
\[
\text{meeting value before }z
=
\gamma \E_z\!\left[V^s_{x,y}(z)\right].
\]
Because $V^s_{x,y}(z)=0$ below the threshold, this expectation integrates only
over accepted draws.

Under Nash bargaining with worker share $\gamma$, the worker's gain from an
accepted meeting is $\gamma V^s_{x,y}(z)$. BS define the value of creating a
meeting before seeing $z$ as the intent-to-treat:
\[
ITT_x(y)
=
\gamma\int_0^\infty V^s_{x,y}(z)g(z)\,dz
=
\gamma\int_{\bar z_{x,y}}^\infty V^s_{x,y}(z)g(z)\,dz.
\]
Source: BS equations (17)--(18), p. 9. The average treatment on the treated,
conditional on the meeting actually becoming a match, is
\[
ATT_x(y)
=
\gamma
\frac{\int_{\bar z_{x,y}}^\infty V^s_{x,y}(z)g(z)\,dz}
{1-G(\bar z_{x,y})}.
\]
Source: BS equation (19), pp. 9--10. Dividing the ITT integral by the ATT
conditional mean gives the derived identity
\[
ITT_x(y)=\left[1-G(\bar z_{x,y})\right]ATT_x(y).
\]
Therefore the comparison is pinned down by the acceptance probability
\[
p_{x,y}\equiv 1-G(\bar z_{x,y}).
\]
If $0<p_{x,y}<1$ and accepted matches have positive surplus, then
\[
0<ITT_x(y)<ATT_x(y).
\]
The intuition is simple: $ATT_x(y)$ conditions on the meeting becoming a match,
while $ITT_x(y)$ values the meeting before knowing whether it will clear the
selection threshold. If $p_{x,y}=1$, then every meeting is accepted and
$ITT_x(y)=ATT_x(y)$. If $p_{x,y}=0$, the ATT is not economically meaningful
because there are no treated/accepted matches.
Under the Pareto case used in BS Section 3, this becomes
\[
ITT_x(y)
=
\left(\frac{\bar z_{x,y}}{z_0}\right)^{-\theta}ATT_x(y).
\]
Source: BS equation (27), pp. 13--14. BS emphasize the empirical content of this
identity: wage data on accepted matches reveal objects tied to the ATT, not the
ITT, unless one also knows the probability that a meeting clears the acceptance
threshold. Source: BS equation (20), p. 10, and Sec. 3.2, pp. 13--14.

\subsection{Is $\pi(x\mid y)$ a Wage?}

For Shimer--Smith, yes in the limited sense that $\pi(x\mid y)$ is the flow payoff
received by type $x$ while matched with type $y$. It is the worker-side flow
payment in the Bellman equation, and feasibility imposes
\[
\pi(x\mid y)+\pi(y\mid x)=f(x,y).
\]
Source: Shimer--Smith, pp. 347--348. The old note derived the useful expression
\[
\pi(x\mid y)
=
\frac12 f(x,y)+\frac12 w(x)-\frac12 w(y).
\tag{SS-payoff}
\]
Source: prior result PDF, Sec. 6.3; \texttt{shimer\_smith\_writeup.tex},
Sec. 6.3. This is not printed in Shimer--Smith; it follows from their equations
(3)--(4), pp. 348--349.

The word ``wage'' needs care. Shimer--Smith are writing a symmetric
transferable-output matching model, not a worker-firm employment model. Thus
$\pi(x\mid y)$ is wage-like because it is the flow payment to the $x$ side while
matched, but the same notation could be used for either side of the match. The
formula says: the $x$ side gets half of the match output, plus half of its own
outside value, minus half of the partner's outside value. Intuitively, a higher
outside value for $x$ raises the payment needed to make $x$ willing to match,
while a higher outside value for $y$ gives $y$ a stronger claim on the same output
pie.

For BS, the actual wage object is not $\pi(x\mid y)$. BS use $\bar\pi_y$ for the
firm's reservation profit, while the accepted-match wage is $W_{x,y}(z)$. The
structural wage equation is
\[
W_{x,y}(z)
=
\bar w_x+\gamma\left(zf_{x,y}-\bar w_x-\bar\pi_y\right),
\qquad z\geq \bar z_{x,y}.
\]
Source: BS equation (9), p. 7. Derived comparison to SS: this is exactly parallel to the SS-payoff formula,
outside value plus worker bargaining share times realized net surplus. Setting
$\gamma=1/2$ and replacing $z f_{x,y}$ by deterministic $f(x,y)$ gives
\[
W_{x,y}
=
\frac12 f_{x,y}+\frac12\bar w_x-\frac12\bar\pi_y,
\]
which is the same additive-plus-half-production structure as SS-payoff.

Thus the clean answer is: $\pi(x\mid y)$ is the wage-like object in
Shimer--Smith, but in BS the wage is $W_{x,y}(z)$, not $\bar\pi_y$. The symbol
$\bar\pi_y$ is the firm's reservation profit, an outside-option object on the
firm side. Confusing $W_{x,y}(z)$ and $\bar\pi_y$ would mix up the worker's actual
pay in an accepted match with the firm's value of remaining vacant.

\subsection{Check Against BS Proposition 2}

Proposition 2 is the payoff from adding the stochastic $z$ and then conditioning
on accepted matches. Under a Pareto distribution for $z$, BS prove that the wage
observed in an accepted match can be written
\[
\omega_{i,j}
=
\alpha^*_{x_i}+\psi^*_{y_j}+\epsilon_{i,j},
\]
where
\[
\alpha^*_x
=
\left(1+\frac{\gamma}{\theta-1}\right)\bar w_x,
\qquad
\psi^*_y
=
\frac{\gamma}{\theta-1}\bar\pi_y,
\]
and the residual is mean zero conditional on the worker and firm types. Source:
BS Proposition 2, pp. 11--12.

BS then connect the firm wage component to treatment effects. In the Pareto
case,
\[
ATT_x(y)
=
\frac{\gamma(\bar w_x+\bar\pi_y)}{(\theta-1)(r+\delta)}
=
\frac{\frac{\gamma}{\gamma+\theta-1}\alpha^*_x+\psi^*_y}{r+\delta}.
\]
Therefore, for a fixed worker type,
\[
ATT_x(y_1)-ATT_x(y_2)
=
\frac{\psi^*_{y_1}-\psi^*_{y_2}}{r+\delta}.
\]
Source: BS equations (24)--(26), p. 13. This is the precise sense in which
firm wage premia line up with the ATT: they recover differences in accepted
match values, not differences in the ex ante value of creating meetings.

Derived comparison: this does line up with the SS additive-plus-$\frac12 f$ result, but with an
important twist. In SS, the payoff surface is additive plus one-half of the
deterministic production surface. In BS with Pareto $z$, conditioning on
acceptance makes the normalized overshoot $z/\bar z_{x,y}$ have the same
distribution for every type pair. Since
\[
\bar z_{x,y}f_{x,y}=\bar w_x+\bar\pi_y,
\]
the accepted wage can be rewritten in terms of $(\bar w_x,\bar\pi_y)$ and the
accepted overshoot, rather than in terms of $f_{x,y}$ directly. That is why
Proposition 2 says the accepted wage is additively separable even though the
matching probabilities still depend on $f_{x,y}$ through $\bar z_{x,y}$.

Here is the Pareto step written out. Start from the BS wage equation:
\[
W_{x,y}(z)
=
\bar w_x+\gamma(zf_{x,y}-\bar w_x-\bar\pi_y).
\]
Use the threshold identity $\bar z_{x,y}f_{x,y}=\bar w_x+\bar\pi_y$ to rewrite
the production term as
\[
zf_{x,y}
=
\frac{z}{\bar z_{x,y}}(\bar w_x+\bar\pi_y).
\]
Now define the accepted-match overshoot
\[
Q=\frac{z}{\bar z_{x,y}}
\qquad\text{conditional on }z\geq \bar z_{x,y}.
\]
Under the Pareto distribution, the conditional distribution of $Q$ is the same for
every type pair:
\[
\Pr(Q\geq q\mid z\geq \bar z_{x,y})
=
\frac{\Pr(z\geq q\bar z_{x,y})}{\Pr(z\geq \bar z_{x,y})}
=
q^{-\theta}.
\]
This Pareto scale-invariance calculation was already written in the prior
stochasticity note. Source: \texttt{model\_acceptance\_stochasticity\_writeup.tex},
Sec. 4; prior result PDF, Sec. 8.4. Because the distribution of $Q$ is common
across type pairs, $f_{x,y}$ disappears from the conditional mean accepted wage
after the threshold substitution. It does not disappear from the model: it still
governs the probability of acceptance, $1-G(\bar z_{x,y})$, and therefore the
match graph and ITT.

So the exact slogan should be:
\[
\text{SS: deterministic type-pair surplus gives payoff } 
\frac12 f+\text{additive outside values.}
\]
\[
\text{BS: stochastic type-pair surplus plus Pareto selection gives accepted wages }
\text{additive in reservation values.}
\]
The wage surface can look additively separable in BS while the match graph is
strongly shaped by nonadditive production through selection thresholds. This is
the real bridge from Bellman to selection to ITT.


% ============================================================
% Merged from bilal_modeling_memo.tex
% ============================================================


\section{Bilal Slides and the Low-Rank Modeling Section}

\paragraph{Source convention.}

Page citations to Bilal refer to the PDF pages in
\texttt{bilal\_2024\_slides\_search-and-matching\_Lecture2.pdf}. Page citations
to Bonhomme--Lamadon--Manresa refer to the printed Econometrica pages in BLM
(2019). Page citations to the project note refer to PDF pages in
\texttt{loo\_full\_note\_v1.pdf}. When I reuse results from the prior project
notes, I cite the prior result PDF \texttt{7\_15\_Meeting (1).pdf} and the
source \TeX{} file \texttt{shimer\_smith\_writeup.tex}. Displays or claims not
directly sourced are labeled as derived synthesis, derived ANOVA steps, or
derived nesting interpretations.

\subsection{Notation}

\begin{longtable}{>{\raggedright\arraybackslash}p{0.24\linewidth}
                >{\raggedright\arraybackslash}p{0.68\linewidth}}
\toprule
Symbol & Meaning \\
\midrule
\endfirsthead
\toprule
Symbol & Meaning \\
\midrule
\endhead
$u_t,v_t$ & Bilal unemployment and vacancy stocks at time $t$. Source: Bilal, pp. 9--10. \\
$M_t,M(u,v)$ & Bilal aggregate matching function and number of matches. Source: Bilal, p. 9. \\
$m,\alpha$ & Bilal Cobb--Douglas matching-function scale and unemployment elasticity in $M(u,v)=mu^\alpha v^{1-\alpha}$. Source: Bilal, p. 9. \\
$\theta_t$ & Bilal labor-market tightness, $v_t/u_t$. It is positive but not necessarily greater than one, since vacancies can be smaller or larger than unemployment. Source: Bilal, p. 10. \\
$f_t,q_t$ & Bilal job-finding and vacancy-filling probabilities. Source: Bilal, p. 10. \\
$s$ & Bilal exogenous job-termination probability. Source: Bilal, p. 10. \\
$U_t,E_t,V_t,J_t$ & Bilal values of unemployment, employment, a vacant job/firm, and a filled job/firm. Source: Bilal, p. 14. \\
$b,c,z_t,w_t,\beta$ & Bilal unemployment flow payoff, vacancy cost, productivity, wage, and discount factor. Source: Bilal, pp. 9, 14. \\
$S_t,\gamma,r$ & Bilal joint surplus, worker bargaining power, and steady-state interest-rate notation $r=\beta^{-1}-1$. Source: Bilal, pp. 15, 19, 23. \\
$Y_{ij},X_{ij}'\delta,\varepsilon_{ij}$ & Project wage/log-wage observation, observed covariate component, and residual in the low-rank interface. Source: project note, p. 8. \\
$i,j,I,J$ & Worker and firm indices and counts in the project interface. Source: project note, p. 8. \\
$\alpha_i,\psi_j,u_i,v_j,\Lambda,\operatorname{rank}(\Lambda)=r$ & Project additive worker and firm components, worker and firm interaction factors, interaction matrix, and interaction rank. Source: project note, p. 8. The symbol $r$ is therefore context-dependent: an interest rate in Bilal's steady-state slide and an interaction rank in the project interface. \\
$\delta,\varepsilon_{ij}$ & Project covariate coefficient vector in $X_{ij}'\delta$ and residual in $(\star)$. Source: project note, p. 8. \\
$k,\ell,K,L$ & Worker type and firm class indices and counts in the saturated type-cell construction. Source: project note, p. 8. \\
$\mu_{k\ell},\mu_{\cdot\cdot},\mu_{k\cdot},\mu_{\cdot\ell}$ & Saturated type-cell mean and its grand, row, and column means. Derived from the ANOVA decomposition in project note Proposition 1, p. 8. \\
$\bar\alpha_k,\bar\psi_\ell,M_{k\ell}$ & Additive type components and double-centered interaction matrix in the ANOVA decomposition. Derived from project note Proposition 1, p. 8. This $M$ is unrelated to Bilal's matching function $M(u,v)$. \\
$e_k,e_\ell,\1_K,\1_L$ & Standard basis vectors and all-ones vectors used in the type-indicator nesting and rank proof. Derived notation for project note Proposition 1, p. 8. \\
$k_{it},\ell(j),\alpha_i,a_t(k),b_t(k),X_{it}'c_t,\varepsilon_{it}$ & BLM firm class, project firm-class map, worker type, class intercept, class slope, covariate component, and residual in the simple static regression example. Source for BLM objects: BLM, pp. 702--704; project notation $\ell(j)$ is derived for the type-cell embedding. \\
\bottomrule
\end{longtable}

\subsection{Bilal Lecture 2: What Matters for This Project}

Bilal's Lecture 2 is a compact DMP search-and-matching benchmark. It is not a
two-sided worker-type/firm-type model like Shimer--Smith or
Borovi\v{c}kov\'a--Shimer, but it gives the reusable search logic: matching
technology, value functions, joint surplus, Nash bargaining, free entry, and
the unemployment stock-flow law.

Derived synthesis for teaching the section: read the slides as a five-step
chain.
\begin{enumerate}
\item A matching function turns unemployment and vacancies into meetings.
\item Tightness summarizes how favorable the market is to workers versus firms.
\item Worker and firm Bellman equations value unemployment, employment,
vacancies, and filled jobs.
\item Adding worker and firm values gives joint surplus, which removes the wage
because the wage is a transfer inside the match.
\item Bargaining splits the surplus, and free entry pins down tightness.
\end{enumerate}
That chain is why the slides matter for this project. They are not a source for
the low-rank wage equation itself; they are the cleanest source for the search
logic that sits behind the SS-to-BS comparison.

\subsubsection{Matching Technology and Tightness}

The lecture defines vacancies $v_t$ and unemployment $u_t$, with matches formed
by a matching function
\[
M_t=M(u_t,v_t).
\]
The slide specializes to Cobb--Douglas,
\[
M(u,v)=m u^\alpha v^{1-\alpha},
\]
while noting that the qualitative results extend to any constant-returns
matching function. Source: Bilal, p. 9.

Labor-market tightness is
\[
\theta_t=\frac{v_t}{u_t}.
\]
The job-finding probability for an unemployed worker is
\[
f_t=\frac{M_t}{u_t}=f(\theta_t)=m\theta_t^{1-\alpha},
\]
and the vacancy-filling probability is
\[
q_t=\frac{M_t}{v_t}=q(\theta_t)=m\theta_t^{-\alpha}
=\frac{f(\theta_t)}{\theta_t}.
\]
Source: Bilal, p. 10. The key conceptual point for our work is that random
matching can be summarized either at the individual meeting-rate level, as in
SS and BS, or at the aggregate matching-function/tightness level, as in DMP.
When $\theta_t$ rises, vacancies are abundant relative to unemployed workers.
Then workers find jobs more easily because $f_t$ rises, while firms fill
vacancies less easily because $q_t$ falls. This is the first place where the
two-sided nature of search appears even though the model has only one aggregate
worker type and one aggregate job type.

\subsubsection{Value Functions}

The lecture's four Bellman equations are
\[
U_t=b+\beta \E_t\left[f_t E_{t+1}+(1-f_t)U_{t+1}\right],
\]
\[
E_t=w_t+\beta \E_t\left[sU_{t+1}+(1-s)E_{t+1}\right],
\]
\[
V_t=-c+\beta \E_t\left[q_tJ_{t+1}+(1-q_t)V_{t+1}\right],
\]
\[
J_t=z_t-w_t+\beta \E_t\left[sV_{t+1}+(1-s)J_{t+1}\right].
\]
Source: Bilal, p. 14. This is the slide-deck analogue of the Bellman logic used
in the SS and BS memo: current payoff plus expected continuation values, with
transitions generated by search and separation.

The equations should be read literally. An unemployed worker receives $b$ today,
then next period is employed with probability $f_t$ and unemployed otherwise. An
employed worker receives $w_t$ today, then loses the job with probability $s$.
A vacancy pays cost $c$ today and is filled with probability $q_t$. A filled job
receives output net of wages, $z_t-w_t$, and separates with probability $s$.
These four equations are the accounting system from which the surplus equation
is obtained.

\subsubsection{Surplus and the Wage}

The worker surplus is $E_t-U_t$, the firm surplus is $J_t-V_t$, and free entry
sets $V_t=0$. The joint match surplus is
\[
S_t=E_t+J_t-U_t,
\]
and the lecture derives
\[
S_t
=
z_t-b-\beta f_t\E_t[E_{t+1}-U_{t+1}]
+\beta(1-s)\E_t[S_{t+1}].
\]
Source: Bilal, p. 15; the negative sign on the foregone-search term is also
confirmed by the later joint-surplus equation on Bilal, p. 21,
$S_t=z_t-b+\beta(1-s-\gamma f_t)\E_t[S_{t+1}]$. The useful point is the same
one that makes SS and BS clean: with transferable utility and linear payoffs,
joint surplus is independent of the wage. The wage only splits the surplus.

The term $z_t-b$ is the current output gain from employment rather than
unemployment. The term $\beta(1-s)\E_t[S_{t+1}]$ is the continuation value of
the match if it survives. The negative foregone-search term appears because a
matched worker cannot simultaneously remain unemployed and keep drawing new job
offers. This nuance is easy to miss: the search opportunity that an unemployed
worker gives up is part of the opportunity cost of being matched.

Bilal then imposes generalized Nash bargaining:
\[
w_t=\arg\max_w (E_t(w)-U_t)^\gamma J_t(w)^{1-\gamma},
\]
so
\[
E_t(w_t)-U_t=\gamma S_t,
\qquad
J_t(w_t)=(1-\gamma)S_t.
\]
Source: Bilal, p. 19. The equilibrium wage can be written as
\[
w_t=(1-\gamma)b+\gamma(z_t+\theta_t c),
\]
after using free entry. Source: Bilal, p. 20. For our project, this is a
warning and a guide: search frictions determine outside options and surplus,
but the wage equation depends on the bargaining/transfer rule. That is why the
SS-to-BS mapping had to track outside values and surplus shares carefully.

\subsubsection{Equilibrium Conditions}

Given surplus sharing, the joint-surplus Bellman becomes
\[
S_t=z_t-b+\beta(1-s-\gamma f_t)\E_t[S_{t+1}],
\]
free entry determines tightness,
\[
c=\beta q_t(1-\gamma)\E_t[S_{t+1}],
\]
and unemployment follows
\[
u_{t+1}-u_t=s(1-u_t)-f_tu_t.
\]
Source: Bilal, p. 21. In steady state, these reduce to the job-creation condition
\[
\frac{r+s+\gamma f(\theta)}{(1-\gamma)q(\theta)}
=
\frac{z-b}{c},
\qquad r=\beta^{-1}-1,
\]
and the Beveridge-curve stock-flow identity
\[
u=\frac{s}{s+f(\theta)}.
\]
Source: Bilal, p. 23.

The equilibrium logic is now closed. Given tightness, the matching function gives
$f(\theta)$ and $q(\theta)$. Given these probabilities, surplus is determined.
Given surplus sharing and free entry, vacancies enter until the expected value
of posting a vacancy equals the cost. The unemployment equation then says how the
stock of unemployed workers evolves from separations into unemployment and job
findings out of unemployment.

\subsubsection{Connection to SS and BS}

Derived synthesis: Bilal's slides should be used as the generic search-model scaffold, not as a
direct nesting target. SS and BS add type-pair heterogeneity to this search
logic. SS makes match output deterministic conditional on types, $f(x,y)$. BS
adds a match-specific draw $z$ and therefore a type-pair selection threshold.
The shared structure is:
\[
\text{Bellman equations}
\quad\Rightarrow\quad
\text{joint surplus independent of wage}
\quad\Rightarrow\quad
\text{Nash split}
\quad\Rightarrow\quad
\text{acceptance/job-creation condition}.
\]

\subsection{Modeling Section: The Low-Rank Interface}

The project note's modeling object is
\[
Y_{ij}
=
X_{ij}'\delta+\alpha_i+\psi_j+u_i^\top\Lambda v_j+\varepsilon_{ij},
\qquad
\operatorname{rank}(\Lambda)=r\ll \min\{I,J\}.
\tag{$\star$}
\]
Source: project note, p. 8; prior result PDF, pp. 18--21;
\texttt{shimer\_smith\_writeup.tex}, Sec. 6. The point of the modeling section is
not to claim that every structural model is literally inside $(\star)$. The point
is to identify which static wage-equation cores land in this interface after
normalization.

\subsubsection{The Double-Demeaned $K\times L$ Construction}

Let worker types be $k\in\{1,\ldots,K\}$ and firm classes be
$\ell\in\{1,\ldots,L\}$. Let
\[
\mu=(\mu_{k\ell})\in\mathbb R^{K\times L}
\]
be the saturated static type-pair mean wage surface:
\[
\mu_{k\ell}
=
\E[Y\mid \text{worker type }k,\text{ firm class }\ell].
\]
This is the object shown in the Saturday screenshot
\texttt{C:\textbackslash Users\textbackslash A8327\textbackslash Downloads\textbackslash image (3).png}
and in the project note's Proposition 1. Source: project note, p. 8.

Derived ANOVA step: define the grand mean, row means, and column means by
\[
\mu_{\cdot\cdot}
=
\frac{1}{KL}\sum_{k=1}^K\sum_{\ell=1}^L\mu_{k\ell},
\qquad
\mu_{k\cdot}
=
\frac{1}{L}\sum_{\ell=1}^L\mu_{k\ell},
\qquad
\mu_{\cdot\ell}
=
\frac{1}{K}\sum_{k=1}^K\mu_{k\ell}.
\]
Set
\[
\bar\alpha_k=\mu_{k\cdot}-\mu_{\cdot\cdot},
\qquad
\bar\psi_\ell=\mu_{\cdot\ell},
\]
and define the double-centered interaction matrix
\[
M_{k\ell}
=
\mu_{k\ell}-\mu_{k\cdot}-\mu_{\cdot\ell}+\mu_{\cdot\cdot}.
\]
Then
\[
\mu_{k\ell}
=
\bar\alpha_k+\bar\psi_\ell+M_{k\ell}.
\]
This is a derived ANOVA decomposition. The steps are:
\begin{enumerate}
\item Start with the fully saturated table $\mu_{k\ell}$, which allows every
worker-type/firm-class cell to have its own mean.
\item Pull out the worker-type level $\bar\alpha_k$ by comparing row $k$ to the
grand mean.
\item Pull out the firm-class level $\bar\psi_\ell$ by using the column mean.
\item Put whatever remains in $M_{k\ell}$.
\end{enumerate}
Because $M_{k\ell}$ is defined as the leftover after removing row and column
levels, it is the pure interaction part of the table. It tells us whether a
particular worker type and firm class have a mean wage above or below what their
separate additive levels would predict. The interaction matrix is double-centered:
\[
\sum_{\ell=1}^L M_{k\ell}=0\quad\text{for every }k,
\qquad
\sum_{k=1}^K M_{k\ell}=0\quad\text{for every }\ell.
\]
The row sums are zero because the average leftover within a row has already been
absorbed into $\bar\alpha_k$. The column sums are zero because the average
leftover within a column has already been absorbed into $\bar\psi_\ell$. It is
unique once the mean convention above is fixed.

Derived embedding in $(\star)$: assign worker $i$ to type $k(i)$ and firm $j$ to class $\ell(j)$. Let
$e_k\in\mathbb R^K$ and $e_\ell\in\mathbb R^L$ be standard basis vectors. Set
\[
\alpha_i=\bar\alpha_{k(i)},
\qquad
\psi_j=\bar\psi_{\ell(j)},
\qquad
u_i=e_{k(i)},
\qquad
v_j=e_{\ell(j)},
\qquad
\Lambda=M.
\]
Then
\[
u_i^\top\Lambda v_j
=
e_{k(i)}^\top M e_{\ell(j)}
=
M_{k(i)\ell(j)},
\]
and therefore
\[
\alpha_i+\psi_j+u_i^\top\Lambda v_j
=
\mu_{k(i)\ell(j)}.
\]
So the saturated $K\times L$ type-pair mean surface is exactly a special case of
$(\star)$. The indicator vectors do no estimation magic here: they simply select
one row and one column of the type-cell matrix. The entire substantive claim is
that the project interface has enough room to reproduce any static type-cell
mean once the additive pieces and the centered interaction matrix are allowed.

\subsubsection{Rank Bound}

Derived rank bound: because $M$ is double-centered, its columns lie in the $K-1$ dimensional
subspace orthogonal to $\1_K$, and its rows lie in the $L-1$ dimensional
subspace orthogonal to $\1_L$. Hence
\[
\operatorname{rank}(M)
\leq
\min\{K-1,L-1\}
=
\min\{K,L\}-1.
\]
This is the rank bound in the screenshot and in the project note. Source:
project note, p. 8. The important interpretation is that the additive row and
column components remove one dimension from each side. The interaction matrix
does not live in all of $\mathbb R^{K\times L}$; it lives between centered type
spaces.

The bound is easiest to explain geometrically. A completely unrestricted
$K\times L$ matrix can vary in $K$ independent row directions and $L$ independent
column directions. Double-centering removes the all-ones direction on the worker
side and the all-ones direction on the firm side. Therefore the interaction can
use at most $K-1$ independent worker directions and at most $L-1$ independent firm
directions, so its rank cannot exceed the smaller of those two numbers.

\subsection{In What Sense We Nest BLM}

\subsubsection{Notation for the BLM Nesting Question}

\begin{longtable}{>{\raggedright\arraybackslash}p{0.24\linewidth}
                >{\raggedright\arraybackslash}p{0.68\linewidth}}
\toprule
Symbol & Meaning \\
\midrule
\endfirsthead
\toprule
Symbol & Meaning \\
\midrule
\endhead
$\mathcal F_\star$ & The class of static conditional mean wage surfaces generated by the project interface $(\star)$. Derived notation for the nesting argument. \\
$\mathcal P_{\mathrm{BLM}}$ & The class of joint panel-data distributions generated by the full BLM framework. Derived notation for the nesting argument. \\
$\Pi_{\mathrm{mean}}$ & Projection that maps a joint panel distribution to its static conditional mean wage surface. Derived notation for the nesting argument. \\
$m_{ij}$ & Static conditional mean wage for worker $i$ at firm $j$, after observed covariates are held fixed or partialled out. Derived notation for $(\star)$. \\
$t,T$ & Calendar period and panel length in BLM. Source: BLM, p. 703. \\
$Y_i^T,j_i^T,k_i^T,m_i^{T-1},X_i^T$ & Worker-level panel histories of earnings, firms, firm classes, mobility indicators, and covariates. Source: BLM, p. 703. \\
$Y_{it},X_{it}$ & Worker $i$'s log earnings and observed covariates in period $t$. Source: BLM, p. 703. \\
$j_{it},k_{it}$ & Firm employing worker $i$ and that firm's latent class in period $t$. Source: BLM, pp. 702--704. \\
$m_{it}$ & BLM mobility indicator between periods. Source: BLM, pp. 703--704. This is unrelated to the derived mean surface $m_{ij}$. \\
$\alpha_i$ & BLM worker type, discrete or continuous depending on the specification. Source: BLM, p. 703. \\
$a_t(k),b_t(k),c_t,\varepsilon_{it}$ & BLM simple static regression intercept, worker-type slope, covariate coefficient, and residual. Source: BLM, p. 704. \\
$F_{t,k,\alpha,x}$ & Generic notation for the conditional distribution of log earnings at time $t$ for worker type $\alpha$, firm class $k$, and covariates $x$. BLM's target is earnings distributions, not only means. Source: BLM, pp. 703--704; notation is derived. \\
$P(k_{i,t+1},m_{it}\mid\cdot),\Pr,\Pr'$ & BLM mobility/transition law and two generic probability laws used in the counterexample. In the dynamic model mobility may depend on current earnings $Y_{it}$. Source: BLM, p. 704; prime notation is derived. \\
$F(Y_{i,t+1}\mid\cdot)$ & BLM post-move or next-period earnings distribution. In the dynamic model it may depend on previous earnings and previous firm class. Source: BLM, p. 704. \\
\bottomrule
\end{longtable}

BLM's static model uses firm classes and worker heterogeneity. Firms are grouped
into classes $k(j)$, and workers have latent types $\alpha_i$, which may be
discrete or continuous. BLM are interested in earnings distributions for workers
of a given type in a given firm class, as well as type proportions and sorting.
Source: BLM, pp. 702--704. A simple static BLM regression example is
\[
Y_{it}=a_t(k_{it})+b_t(k_{it})\alpha_i+X_{it}'c_t+\varepsilon_{it}.
\]
Source: BLM, p. 704.

Derived nesting interpretation: the honest nesting claim has three layers.

\paragraph{1. Exact for static type-cell mean surfaces.}
If BLM worker heterogeneity is discretized into $K$ worker types and firm
heterogeneity into $L$ firm classes, then the static conditional mean surface
\[
\mu_{k\ell}
=
\E[Y\mid k(i)=k,\ell(j)=\ell]
\]
is exactly nested by $(\star)$ through the double-demeaned construction above.
This is the clean reading of the screenshot's BLM row.
The word ``static'' is doing real work. We are only talking about the mean wage
attached to a worker type and firm class at a point in the model, after the
types/classes have been defined.

\paragraph{2. Exact for BLM's simple interactive mean example.}
The displayed BLM regression
\[
a_t(k)+b_t(k)\alpha_i
\]
is a firm-class intercept plus a firm-class slope on worker type. For a fixed
date $t$, it is nested by
\[
\psi_j=a_t(\ell(j)),
\qquad
u_i=\alpha_i,
\qquad
v_j=b_t(\ell(j)),
\qquad
\Lambda=1,
\]
with observed covariates absorbed into $X_{ij}'\delta$. This is rank one if
$\alpha_i$ is scalar. If the empirical object is instead a fully saturated
discrete type-class mean matrix, the exact representation is the $M$ construction
above, with rank at most $\min\{K,L\}-1$.

This is also where the two meanings of ``low rank'' need to be kept separate.
The simple BLM mean example is algebraically rank one because the interaction is
the product of a worker scalar $\alpha_i$ and a firm-class slope $b_t(k)$. But it
does not necessarily reduce the effective number of worker parameters: every
worker can still carry their own $\alpha_i$. The project note's nesting claim is
about algebraic rank of the interaction surface, not automatically about a small
number of worker parameters.

\paragraph{3. Not exact for the full BLM framework.}
Here is the rigorous version of the non-nesting claim. The project interface is
a model for a static conditional mean surface. After suppressing covariates for
notation, it says
\[
m_{ij}
\equiv
\E[Y_{ij}\mid i,j]
=
\alpha_i+\psi_j+u_i^\top\Lambda v_j.
\tag{LR-mean}
\]
Thus $\mathcal F_\star$ is a set of maps
\[
(i,j)\mapsto m_{ij}.
\]
It does not specify a joint distribution for the worker panel
\[
\left(Y_i^T,j_i^T,k_i^T,m_i^{T-1},X_i^T\right),
\]
nor does it specify a law for firm-class assignment, worker-type proportions,
mobility, or the full conditional distribution of earnings around the mean.

BLM's full framework is a model for those additional objects. The observed data
are sequences of earnings, firm and mobility indicators, and covariates. Source:
BLM, p. 703. BLM's target includes earnings distributions for workers of type
$\alpha$ in firms of class $k$ and the proportions of worker types in firm
classes. Source: BLM, p. 703. Their static example includes the mean regression
\[
Y_{it}=a_t(k_{it})+b_t(k_{it})\alpha_i+X_{it}'c_t+\varepsilon_{it},
\]
but the framework around it is distributional and panel-based. Source: BLM,
p. 704. Their dynamic model additionally allows current earnings to affect
mobility and allows earnings in period $t+1$ to depend on previous earnings and
previous firm class. Source: BLM, p. 704.

Therefore, if ``model A nests model B'' means that every observable distribution
generated by B can be generated by A under some parameter choice, then
$(\star)$ does not nest the full BLM framework. The reason is not just verbal; it
is a domain mismatch:
\[
\mathcal F_\star
\subseteq
\left\{\text{static conditional mean surfaces }(i,j)\mapsto m_{ij}\right\},
\]
whereas full BLM lives in a space like
\[
\mathcal P_{\mathrm{BLM}}
\subseteq
\left\{
\text{joint laws of }(Y_i^T,j_i^T,k_i^T,m_i^{T-1},X_i^T,\alpha_i)
\right\}.
\]
There is no parameter vector $(\delta,\alpha_i,\psi_j,u_i,\Lambda,v_j)$ in
$(\star)$ that determines the BLM objects
\[
F_{t,k,\alpha,x},\qquad
\Pr(\alpha_i\mid k_{it}=k),\qquad
\Pr(k_{i,t+1},m_{it}\mid Y_{it},\alpha_i,k_{it},X_{it}),
\]
or
\[
F(Y_{i,t+1}\mid Y_{it},\alpha_i,k_{it},k_{i,t+1},X_{it},X_{i,t+1},m_{it}).
\]
Those objects are part of BLM's empirical framework; they are not coordinates of
the low-rank wage surface.

A concrete counterexample makes the point. Take two BLM data-generating processes
with the same static mean
\[
\E[Y_{it}\mid \alpha_i,k_{it},X_{it}]
=
a_t(k_{it})+b_t(k_{it})\alpha_i+X_{it}'c_t,
\]
but different conditional earnings distributions, for example one with small
conditional variance and one with large conditional variance. Both have the same
mean surface and therefore the same representation in $(\star)$, but they imply
different earnings distributions. Since BLM's target includes the earnings
distribution, $(\star)$ cannot be a nesting of the full BLM model.

Likewise, take two BLM dynamic models with the same contemporaneous mean surface
but different mobility laws:
\[
\Pr(m_{it}=1\mid Y_{it},\alpha_i,k_{it},X_{it})
\neq
\Pr'(m_{it}=1\mid Y_{it},\alpha_i,k_{it},X_{it}).
\]
They again map to the same static low-rank mean surface, but they imply different
joint panel distributions and different sorting/mobility patterns. Since
$(\star)$ has no transition law for $m_{it}$ or $k_{i,t+1}$, it cannot reproduce
or restrict those differences. Adding an arbitrary mobility model on top of
$(\star)$ would be a new augmented model, not the original low-rank interface.

So the exact mathematical statement is about a projection, not a full nesting.
Let $\Pi_{\mathrm{mean}}$ map a BLM joint panel distribution to the static
conditional mean surface it implies. For the static mean layers discussed above,
\[
\Pi_{\mathrm{mean}}\left(\mathcal P_{\mathrm{BLM}}\right)
\subseteq \mathcal F_\star
\quad\text{under the embeddings above,}
\]
when attention is restricted to the discretized type-cell means or the simple
interactive mean example. But this does not imply
\[
\mathcal P_{\mathrm{BLM}}\subseteq \mathcal F_\star.
\]
Indeed, the latter expression is not even a well-typed nesting claim: elements
of $\mathcal F_\star$ are mean functions, while elements of
$\mathcal P_{\mathrm{BLM}}$ are joint distributions. The projection
$\Pi_{\mathrm{mean}}$ is many-to-one. It deliberately discards distributions,
type proportions, class-estimation/mixture structure, and dynamic transition
laws. BLM themselves emphasize that the dynamic model allows earnings to affect
mobility and allows post-move earnings to depend on previous earnings and
previous firm class. Source: BLM, p. 704. They also note that still richer
non-Markovian structures can fall outside their own framework. Source: BLM,
p. 706. This reinforces the point: BLM is a framework for panel distributions
and dynamics, not merely a single static wage-surface equation.

So the phrase ``we nest BLM'' is correct only if it is expanded as:
\[
\text{we nest BLM's static wage-equation/type-cell mean layer.}
\]
It should not be written as:
\[
\text{we nest the full BLM empirical framework.}
\]
This matches the prior result PDF's warning that the factor interface nests
static wage-equation cores, while BLM's grouping, mixture estimation, and dynamic
extensions remain separate implementation and interpretation layers. Source:
prior result PDF, pp. 18--21; \texttt{shimer\_smith\_writeup.tex}, Sec. 6.5.

Put differently, the surface nests more than the factors. If two different
normalizations or latent-type representations generate the same static mean
surface, $(\star)$ nests the observational surface. It does not by itself identify
which worker factor, firm factor, or class assignment was the structural one.

\subsection{Clean Wording for the Modeling Section}

The safest paragraph is:

\begin{quote}
The low-rank interface nests the static wage-equation cores of several models.
For BLM, the exact statement is about the static type-cell mean surface. If
worker types are $k=1,\ldots,K$ and firm classes are $\ell=1,\ldots,L$, any
saturated mean matrix $\mu_{k\ell}$ decomposes uniquely into additive row and
column terms plus a double-centered interaction matrix $M$. Taking worker and
firm factors to be type indicators and $\Lambda=M$ reproduces the mean surface
exactly, with $\operatorname{rank}(M)\leq\min\{K,L\}-1$. This nests BLM's static
mean layer, but not BLM's full distributional estimator, class-estimation step,
or dynamic panel model. Mathematically, the low-rank interface is a family of
static conditional mean surfaces, while full BLM is a family of joint panel
distributions with earnings distributions, worker-type proportions, firm-class
assignment, and mobility laws. The former can represent the mean layer of the
latter, but it does not determine or nest the latter as a full empirical
framework.
\end{quote}


% ============================================================
% Merged from misc_eventstudy_kline_memo.tex
% ============================================================


\section{Sorkin--Warwar Event Studies and Kline Edge Effects}

\paragraph{Source convention.}

Page citations to Sorkin--Warwar refer to the PDF pages in
\texttt{sorkin warwar 2026 sw\_logadditivity.pdf}. Page citations to Kline refer
to the PDF pages in \texttt{kline\_2024\_review\_akm\_firm\_wage\_effects\_w33084.pdf}.
When I reuse the prior project write-up, I cite the prior result PDF
\texttt{7\_15\_Meeting (1).pdf} and \texttt{shimer\_smith\_writeup.tex}.
Displays or claims not directly sourced are labeled as derived examples,
derived algebra, or derived mappings.

\subsection{Notation}

\begin{longtable}{>{\raggedright\arraybackslash}p{0.24\linewidth}
                >{\raggedright\arraybackslash}p{0.68\linewidth}}
\toprule
Symbol & Meaning \\
\midrule
\endfirsthead
\toprule
Symbol & Meaning \\
\midrule
\endhead
$Y_{it},Y_{it}(j)$ & Observed wage and potential wage for worker $i$ at time $t$ if assigned to firm $j$. Source for Kline potential-outcome notation: Kline, pp. 29--31. \\
$D_{it}$ & Firm employing worker $i$ at time $t$ in Kline's causal notation. Source: Kline, p. 30. \\
$\alpha_i,\psi_j,X_{it}'\beta,\varepsilon_{it},\mu_{ij}$ & Worker effect, firm effect, covariate component, AKM error, and illustrative time-invariant worker-firm match effect. Source for AKM components: Kline, p. 15. The match-effect extension is derived for the separability discussion. \\
$\tau_t,\varepsilon_{it}(j)$ & Illustrative common time component and transitory potential-wage shock used in the derived parallel-trends example. \\
$A,B,a,b$ & Origin and destination firms and two paired movers in the derived paired-mover algebra. Source for paired-mover design: SW abstract, p. 1. \\
$G_A,G_B,\beta_{\mathrm{paired}}$ & Origin wage gap, destination wage gap, and paired-mover forecast coefficient. Derived notation for the paired-mover separability calculation. \\
$h_i,p_j,m_{ij}$ & Worker, firm, and match effects in SW's partial-equilibrium search model. Source: SW, p. 27. \\
$EE,EUE,\lambda_1$ & Employer-to-employer moves, employer-unemployment-employer moves, and SW's on-the-job offer-arrival parameter. Source: SW, pp. 27, 31--33. \\
$X,Y,g,\beta_{\mathrm{pooled}}$ & Generic regressor, outcome, group, and pooled OLS coefficient in the derived pooling identity. \\
$R,E,B,\Delta,u,\hat\Delta,\tilde\Delta$ & Kline adjusted wage changes, worker-by-edge matrix, firm-by-edge incidence matrix, unrestricted directed edge effects, edge-model error, estimated edge effects, and AKM-predicted edge effects. Source: Kline, pp. 18--19, 24--27. \\
$c,C,\dot\psi,\dot\eta$ & Kline cycle vector, matrix collecting fundamental cycles, projected firm-effect coefficients, and cycle-effect coefficients. Source: Kline, pp. 18--20. \\
$W,H,M,h_{\ell\ell},n_\ell,\ell$ & Kline mover-weight matrix, projection matrix, residual-maker matrix, edge leverage, number of movers on edge $\ell$, and edge index. Source: Kline, pp. 24--27. \\
$Y_{ij},X_{ij}'\delta,u_i,v_j,\Lambda$ & Project low-rank wage-interface notation used in the nesting comparison. Source: prior result PDF, pp. 18--21; \texttt{shimer\_smith\_writeup.tex}, Sec. 6. \\
$m_{ij},\mu_{jk}$ & Derived static mean wage surface and edge-specific mean worker factor, $\mu_{jk}=\E[u_i\mid i:j\to k]$, used to map low-rank surfaces into Kline edge effects. \\
\bottomrule
\end{longtable}

\subsection{Sorkin--Warwar: Why the Event Study Can Have Low Power}

This section is self-contained. Its purpose is to explain one narrow question:
why can a standard AKM event-study test look reassuring even when worker-firm
match effects are important?

\subsubsection{Local Notation}

\begin{longtable}{>{\raggedright\arraybackslash}p{0.24\linewidth}
                >{\raggedright\arraybackslash}p{0.68\linewidth}}
\toprule
Symbol & Meaning \\
\midrule
\endfirsthead
\toprule
Symbol & Meaning \\
\midrule
\endhead
$i$ & Worker index. \\
$A,B,j,k$ & Firms. $A$ and $B$ are used for a common origin-destination pair in the paired-mover calculation. \\
$Y_{it}$ & Observed log earnings of worker $i$ at time $t$. Source: SW, p. 31, for changes in log earnings in the event-study exercise. \\
$Y_i(j)$ & Static potential log earnings of worker $i$ at firm $j$ in the derived separability calculation. \\
$\alpha_i,\psi_j$ & Worker and firm components in the additive AKM wage model. Source for AKM additive structure: Kline, p. 15; used here as the benchmark tested by SW. \\
$\mu_{ij}$ & Worker-firm match effect in the derived wage decomposition. This is the object ruled out by additive separability. \\
$h_i,p_j,m_{ij}$ & Worker, firm, and match payoff components in SW's model. Source: SW, p. 27. \\
$\Delta Y_i$ & Realized earnings change for a mover, destination earnings minus origin earnings. Source: SW, pp. 31--33. \\
$\Delta\psi_i$ & Change in firm effects for a mover, destination firm effect minus origin firm effect. Source: SW, pp. 31--33. \\
$\beta_{\mathrm{ES}}$ & Event-study regression coefficient of $\Delta Y_i$ on $\Delta\psi_i$. Source: SW, pp. 31--33. \\
$G_A,G_B$ & Pay gap between two paired movers at the origin firm $A$ and destination firm $B$. Source for paired-mover design: SW abstract, p. 1. \\
$\beta_{\mathrm{paired}}$ & Forecast coefficient from destination pay gaps on origin pay gaps in the paired-mover design. Source: SW abstract, p. 1. \\
$EE,EUE$ & Employer-to-employer moves and employer-unemployment-employer moves. Source: SW, pp. 31--33. \\
$g,X,Y$ & Generic group, regressor, and outcome used in the derived pooled-slope identity. \\
\bottomrule
\end{longtable}

\subsubsection{What the Event Study Tests}

SW describe the event-study test as asking whether changes in firm effects, or
mean co-worker earnings, predict within-worker changes in earnings. Source: SW,
p. 31. In the notation here, the test is the slope in
\[
\Delta Y_i
=
\beta_{\mathrm{ES}}\Delta\psi_i+\text{error}_i.
\]
SW construct it by randomly splitting workers, estimating two sets of firm
effects, and instrumenting the change in firm effects estimated in one sample
with the change estimated in the other sample. Source: SW, p. 32. The split-sample
instrument is meant to address measurement error in the generated firm effects.

Under the additive AKM benchmark, the wage of worker $i$ at firm $j$ is
\[
Y_i(j)=\alpha_i+\psi_j.
\tag{AKM}
\]
If worker $i$ moves from firm $j$ to firm $k$, the worker component cancels:
\[
\Delta Y_i
=
Y_i(k)-Y_i(j)
=
\psi_k-\psi_j
=
\Delta\psi_i.
\]
So in the additive benchmark, the event-study coefficient should be
\[
\beta_{\mathrm{ES}}=1.
\]
This is the reason the event-study is a useful diagnostic: if workers moving to
higher firm-effect firms do not receive correspondingly higher earnings, additive
firm effects or exogenous mobility look suspicious.

\subsubsection{Why This Is Not the Same as Testing Separability}

Additive separability rules out worker-firm match effects. To see the distinction,
allow the static wage at firm $j$ to be
\[
Y_i(j)=\alpha_i+\psi_j+\mu_{ij},
\tag{Match}
\]
where $\mu_{ij}$ is a worker-firm match effect. For a mover from $j$ to $k$,
\[
\Delta Y_i
=
\psi_k-\psi_j+\mu_{ik}-\mu_{ij}
=
\Delta\psi_i+\Delta\mu_i.
\]
The population event-study slope is therefore
\[
\beta_{\mathrm{ES}}
=
\frac{\Cov(\Delta Y_i,\Delta\psi_i)}{\Var(\Delta\psi_i)}
=
1+
\frac{\Cov(\Delta\mu_i,\Delta\psi_i)}{\Var(\Delta\psi_i)}.
\tag{ES-slope}
\]
This formula is the key. The event-study coefficient can equal one even when
match effects are large, as long as the change in match quality is uncorrelated
with the change in firm effects in the sample used by the regression:
\[
\Cov(\Delta\mu_i,\Delta\psi_i)=0
\quad\Rightarrow\quad
\beta_{\mathrm{ES}}=1,
\]
even if
\[
\Var(\Delta\mu_i)>0.
\]
Thus the event study is not a direct test of $\mu_{ij}=0$. It tests whether the
match-effect change is systematically aligned with the firm-effect change.

The same point can be stated using pre-trends. If $\mu_{ij}$ is time-invariant,
then it sits in the level of the worker-firm wage and need not generate a
pre-trend. A worker can have a special match value at firm $j$ that is constant
over time. That match value violates additive separability, but by itself it does
not imply that the worker's wage would have been trending differently before the
move. This is why passing a pre-trend or event-study diagnostic is weaker than
showing that wages are additively separable.

\subsubsection{Why Paired Movers Test Separability More Directly}

SW's paired-mover design asks a different question. Consider two workers, $a$ and
$b$, who both move from firm $A$ to firm $B$. Their wage gap at the origin is
\[
G_A
=
Y_a(A)-Y_b(A)
=
(\alpha_a-\alpha_b)+(\mu_{aA}-\mu_{bA}),
\]
and their wage gap at the destination is
\[
G_B
=
Y_a(B)-Y_b(B)
=
(\alpha_a-\alpha_b)+(\mu_{aB}-\mu_{bB}).
\]
The firm component cancels because both workers are observed at the same origin
firm and the same destination firm. Under additive separability, $\mu_{ij}=0$ for
all $(i,j)$, so
\[
G_B=G_A.
\]
That is the paired-mover implication: the pay gap between the two workers should
be portable from the origin firm to the destination firm.

With match effects, the gap need not be portable. In the simple derived case
where match effects are independent across firms and independent of worker
effects,
\[
\beta_{\mathrm{paired}}
=
\frac{\Cov(G_B,G_A)}{\Var(G_A)}
=
\frac{\Var(\alpha_a-\alpha_b)}
{\Var(\alpha_a-\alpha_b)+\Var(\mu_{aA}-\mu_{bA})}
<1.
\]
SW's abstract states the design in exactly these terms: for two workers both
moving from firm $A$ to firm $B$, additive worker and firm effects imply that the
origin pay gap should predict the destination pay gap with coefficient one.
Source: SW abstract, p. 1. They estimate coefficients of $0.73$ in Brazil and
$0.78$ in Italy, which they interpret as evidence of match effects. Source: SW
abstract, p. 1.

\subsubsection{The SW Low-Power Mechanism}

SW then ask whether a model with match effects and endogenous mobility can still
pass the standard event-study test. Their model gives workers payoff
\[
h_i+p_j+m_{ij},
\]
where $h_i$ is a worker effect, $p_j$ is a firm effect, and $m_{ij}$ is a match
effect. Workers observe both firm and match components when deciding whether to
accept offers. Source: SW, p. 27.

In simulated data, the event-study test does detect the problem in one important
subsample. For employer-to-employer moves, where workers have more scope to
select on match quality, SW report an event-study coefficient of $0.57$. For
employer-unemployment-employer moves, they report $0.99$. Source: SW, pp.
31--33. The surprising result is the pooled coefficient: even though $41\%$ of
moves are employer-to-employer moves, the full-sample coefficient is $0.99$.
Source: SW, pp. 31--33. So the pooled event study passes even though a large
subsample fails.

The reason is that a pooled regression slope is not just an average of subgroup
slopes. If the sample is split into groups $g$, the pooled slope of outcome $Y$
on regressor $X$ is
\[
\beta_{\mathrm{pooled}}
=
\frac{\sum_g \Pr(g)\Cov(X,Y\mid g)
+\Cov(\E[X\mid g],\E[Y\mid g])}
{\sum_g \Pr(g)\Var(X\mid g)+\Var(\E[X\mid g])}.
\tag{pooled}
\]
The first term in the numerator is the within-group contribution. The second is
the between-group contribution. SW emphasize this exact logic: pooled
coefficients combine relative sample sizes, within-group variances, and
between-group differences in sample means. Source: SW, pp. 32--33. If the
between-group slope is above one, it can pull the pooled event-study coefficient
back toward one even when the within-group coefficients reveal failures.

SW show this in their simulations. As the share of employer-to-employer moves
rises, both the EE and EUE event-study coefficients fall, but the pooled
coefficient remains above $0.90$ and can exceed either subgroup coefficient
because the between-sample slope is above one. Source: SW, pp. 32--34. In the
Brazilian data, they report an event-study coefficient of $1.09$; in Veneto, they
report $1.18$. Source: SW, p. 34. They also show that splitting the Brazilian
data by whether firm effects increase or decrease gives less reassuring subgroup
coefficients than the pooled regression. Source: SW, p. 34.

\subsubsection{Self-Contained Bottom Line}

The event-study test has low power against match effects for two distinct
reasons.

\begin{enumerate}
\item Algebraically, the event-study slope is
\[
\beta_{\mathrm{ES}}
=
1+\frac{\Cov(\Delta\mu_i,\Delta\psi_i)}{\Var(\Delta\psi_i)}.
\]
It can equal one even when match effects matter, provided changes in match
quality are not correlated with changes in firm effects in the tested sample.

\item Empirically, even when some subgroups fail the event-study test, pooling can
hide the failure because the pooled slope includes between-group mean
differences. This is the specific low-power mechanism SW document on pp. 31--35.
\end{enumerate}

Paired movers are sharper for additive separability because they compare two
workers on the same $A\to B$ transition. The firm movement is common to both
workers, so the test asks whether worker wage gaps are portable across firms. That
portability implication is exactly what worker-firm match effects break.

\subsection{Kline Section 3: What the Content Is}

Kline Section 3 rewrites AKM as a restriction on directed mover edge effects. Start
from
\[
Y_{it}=\alpha_i+\psi_{j(i,t)}+X_{it}'\beta+\varepsilon_{it}.
\]
Source: Kline, p. 15. For two-period movers, define adjusted wage changes
\[
R=Y_2-Y_1-(X_2-X_1)'\beta.
\]
Let $E$ be the worker-by-edge matrix of directed origin-destination indicators,
and let $B$ be the firm-by-edge incidence matrix. Kline shows that AKM implies
\[
R=EB'\psi+\varepsilon.
\]
Source: Kline, pp. 18--19.

The unrestricted edge-effect model is
\[
R=E\Delta+u,
\]
where $\Delta_{jk}$ is the average adjusted wage change for movers from firm $j$
to firm $k$. Source: Kline, p. 19. AKM imposes
\[
\Delta=B'\psi.
\]
This reduces potentially $|E|$ directed edge effects to $J-1$ independent firm
effects. Source: Kline, p. 19.

The easiest way to say the content is: Kline starts with wage changes on directed
firm-to-firm edges and asks when those edge changes can be summarized by firm
levels. If the wage gain from $j$ to $k$ is always $\psi_k-\psi_j$, then moving
around any closed loop must add up to zero. If a loop has a nonzero sum, then the
edge effects cannot all be differences of one global firm-effect vector.

\subsubsection{The Cycle Content}

For any cycle vector $c$ in the mobility graph, $Bc=0$. Therefore AKM implies
\[
c'\Delta=c'B'\psi=0.
\]
Source: Kline, pp. 18--19. Conversely, within a connected set, these cycle
restrictions exhaust the restrictions that AKM imposes on edge effects. Kline
writes
\[
\Delta=B'\dot\psi+C\dot\eta,
\]
where $C\dot\eta$ is the cycle component, and AKM is the restriction
$\dot\eta=0$. Source: Kline, p. 20.

This is the actual content of Section 3: firm effects are not primitive magic
objects. They are a low-dimensional potential function for directed mover wage
changes. They are valid as global firm wage levels only when edge effects are
path-independent around cycles.

\subsubsection{The Noise-Adjusted Goodness-of-Fit Content}

Kline also asks how well AKM approximates unrestricted edge effects in Veneto.
The key object is the residual between estimated unrestricted edge effects and
AKM-predicted edge effects:
\[
\hat\Delta-\tilde\Delta=M\hat\Delta,
\qquad M=I-H.
\]
For the mover-weighted residual sum of squares, Kline decomposes
\[
\E_u[\hat\Delta'M'WM\hat\Delta]
=
\Delta'M'WM\Delta
+\operatorname{trace}\!\left(M'WMV_u[\hat\Delta]\right).
\]
Source: Kline, p. 27. The first term is model error; the second is sampling noise.
Under exact AKM, the first term vanishes but the second remains.

Empirically, among non-bridge edges with at least two movers, the residual sum of
squares is $360.51$, while expected noise under AKM is $207.58$. The difference
$152.92$ is the estimated squared approximation error. Relative to the
noise-adjusted edge-effect sum of squares $978.67$, AKM captures about $84\%$ of
true edge-effect variation; the implied correlation is $0.92$. Source: Kline,
pp. 27--28. Across edge samples, Kline reports noise-adjusted $R^2$ values roughly
between $72\%$ and $88\%$, and about $81\%$ when bridges are included. Source:
Kline, pp. 28--29.

\subsubsection{The Causal Content}

Kline's causal section says that edge effects can be causal under three conditions:
exclusion, parallel trends, and stationarity. Source: Kline, pp. 29--31. Proposition
1 states that, under those assumptions,
\[
\E[Y_{i2}-Y_{i1}\mid D_{i1}=j,D_{i2}=k]=\Delta_{jk},
\]
the average treatment effect on movers along edge $j\to k$. Source: Kline, p. 31.
Crucially, this does not require additive separability. Kline explicitly emphasizes
that edge effects can be causal even when firm effects are heterogeneous, but those
edge effects may be intransitive and therefore may not aggregate to a single global
firm ranking. Source: Kline, pp. 30--31.

\subsection{Is Kline's General Edge Model Nested in Our Low-Rank Factor Model?}

Short answer:
\[
\text{AKM is nested; unrestricted Kline edge effects are not.}
\]

\subsubsection{Local Notation}

\begin{longtable}{>{\raggedright\arraybackslash}p{0.24\linewidth}
                >{\raggedright\arraybackslash}p{0.68\linewidth}}
\toprule
Symbol & Meaning \\
\midrule
\endfirsthead
\toprule
Symbol & Meaning \\
\midrule
\endhead
$i$ & Worker index. \\
$j,k$ & Origin and destination firms. \\
$m_{ij}$ & Static conditional mean wage of worker $i$ at firm $j$, after covariates are suppressed. Derived notation for the low-rank surface. \\
$\alpha_i,\psi_j$ & Worker and firm additive components in the wage surface. Source for AKM components: Kline, p. 15. \\
$u_i,v_j,\Lambda$ & Worker factor, firm factor, and low-rank interaction matrix in the project interface. Source: prior result PDF, pp. 18--21; \texttt{shimer\_smith\_writeup.tex}, Sec. 6. \\
$d_u,d_v$ & Dimensions of the worker and firm factor vectors in the local dimension-count argument. Derived notation. \\
$R,E,B,\Delta,u$ & Kline adjusted wage changes, worker-by-edge matrix, firm-by-edge incidence matrix, unrestricted directed edge effects, and residual. Source: Kline, pp. 18--19. \\
$\Delta_{jk}$ & Kline unrestricted average adjusted wage change for movers from firm $j$ to firm $k$. Source: Kline, p. 19. \\
$\mu_{jk}$ & Mean worker factor among movers on edge $j\to k$, $\mu_{jk}=\E[u_i\mid i:j\to k]$. Derived notation for mapping low-rank wage surfaces into Kline edge effects. \\
$q_{jk}$ & Nonadditive part of the edge effect after removing the firm-effect difference, $q_{jk}=\Delta_{jk}-(\psi_k-\psi_j)$. Derived notation. \\
$C$ & Directed cycle in the firm mobility graph. Source for cycle logic: Kline, pp. 18--20. \\
\bottomrule
\end{longtable}

\subsubsection{Step 1: The Two Models Live on Different Objects}

The project low-rank model is a model of wage levels. Suppressing covariates, its
static mean surface is
\[
m_{ij}
=
\alpha_i+\psi_j+u_i^\top\Lambda v_j.
\tag{LR-level}
\]
This object answers the question: what is worker $i$'s mean wage if employed at
firm $j$?

Kline's unrestricted edge model is a model of mover wage changes. For two-period
movers, Kline defines adjusted wage changes $R$ and writes the unrestricted edge
model as
\[
R=E\Delta+u.
\tag{Kline-edge}
\]
Source: Kline, p. 19. Here $\Delta_{jk}$ is a free average adjusted wage change
for movers from origin firm $j$ to destination firm $k$. This object answers a
different question: among the selected workers who move from $j$ to $k$, what is
their average wage change?

So the logical direction is:
\[
\text{wage-level model}
\quad\Longrightarrow\quad
\text{implied mover edge effects}.
\]
The reverse direction is not automatic. Knowing the average wage change on each
observed edge does not, by itself, recover a worker-by-firm wage-level surface.

\subsubsection{Step 2: AKM Is the Easy Nested Case}

AKM is nested because it is the rank-zero case of the project interface. Set
\[
\Lambda=0.
\]
Then (LR-level) becomes
\[
m_{ij}=\alpha_i+\psi_j.
\tag{AKM-level}
\]
For a worker moving from firm $j$ to firm $k$,
\[
m_{ik}-m_{ij}
=
(\alpha_i+\psi_k)-(\alpha_i+\psi_j)
=
\psi_k-\psi_j.
\]
Therefore AKM implies
\[
\Delta_{jk}=\psi_k-\psi_j.
\tag{AKM-edge}
\]
This is exactly Kline's AKM restriction $\Delta=B'\psi$: edge effects must be
differences of a single vector of firm effects. Source: Kline, pp. 18--19.
Equivalently, the sum of AKM edge effects around any directed cycle is zero,
because the firm effects telescope:
\[
(\psi_B-\psi_A)+(\psi_C-\psi_B)+(\psi_A-\psi_C)=0.
\]
Source for the cycle interpretation: Kline, pp. 18--20.

\subsubsection{Step 3: A Low-Rank Wage Surface Implies Structured Edge Effects}

Now allow the low-rank interaction. From (LR-level), a worker moving from $j$ to
$k$ has wage change
\[
m_{ik}-m_{ij}
=
\psi_k-\psi_j+u_i^\top\Lambda(v_k-v_j).
\]
Averaging over the selected movers on edge $j\to k$ gives
\[
\Delta_{jk}
=
\psi_k-\psi_j+\mu_{jk}^\top\Lambda(v_k-v_j),
\qquad
\mu_{jk}\equiv \E[u_i\mid i:j\to k].
\tag{LR-edge}
\]
This is the clean bridge from the low-rank wage surface to Kline's edge language.
The first term, $\psi_k-\psi_j$, is the AKM part. The second term is the
interaction-composition part. It depends on three objects at once:
\[
\text{mover composition } \mu_{jk},
\qquad
\text{firm-factor difference } v_k-v_j,
\qquad
\text{interaction matrix } \Lambda.
\]

Around a directed firm cycle $C$, the additive firm effects still telescope:
\[
\sum_{(j,k)\in C}(\psi_k-\psi_j)=0.
\]
Thus any cycle sum under the low-rank model must come from the interaction term:
\[
\sum_{(j,k)\in C}\Delta_{jk}
=
\sum_{(j,k)\in C}\mu_{jk}^\top\Lambda(v_k-v_j).
\tag{LR-cycle}
\]
This equation is useful because it shows both sides of the answer. Low-rank wage
surfaces can generate nonzero Kline cycle components, so they are more flexible
than AKM. But the cycle components are not arbitrary: they must be generated by
the same worker factors, firm factors, and low-rank matrix that generate every
other edge.

\subsubsection{Step 4: Kline's Unrestricted Edge Model Is Not Nested}

Kline's unrestricted model places a free parameter on each observed directed
edge:
\[
\Delta=(\Delta_{jk})_{j\to k}.
\]
By contrast, the low-rank model requires every edge to satisfy the representation
\[
\Delta_{jk}
=
\psi_k-\psi_j+\mu_{jk}^\top\Lambda(v_k-v_j).
\]
Equivalently, after removing the additive firm-effect difference,
\[
q_{jk}
\equiv
\Delta_{jk}-(\psi_k-\psi_j)
\]
must satisfy
\[
q_{jk}=\mu_{jk}^\top\Lambda(v_k-v_j)
\quad\text{for every observed edge }j\to k.
\tag{restriction}
\]
This is a genuine restriction. The same $\Lambda$ and the same firm factors
$v_j$ must rationalize all edges at once. An unrestricted Kline edge model does
not impose this shared-factor structure.

A simple special case makes the non-nesting immediate. If $\Lambda=0$, then
$q_{jk}=0$ for every edge, so every directed cycle must sum to zero. But Kline's
unrestricted $\Delta$ allows a three-firm cycle such as
\[
\Delta_{AB}=1,\qquad \Delta_{BC}=1,\qquad \Delta_{CA}=1,
\]
whose cycle sum is $3$. This edge table is allowed by the unrestricted Kline
model but impossible under AKM. Source for unrestricted edge effects and cycle
restrictions: Kline, pp. 18--20.

With $\Lambda\neq 0$, the low-rank model allows nonzero cycle sums, but only of
the form (LR-cycle). For fixed low factor dimensions and fixed mover-composition
moments $\mu_{jk}$, the set of edge residuals
\[
(q_{jk})_{j\to k}
\]
that can be written as $\mu_{jk}^\top\Lambda(v_k-v_j)$ is a structured
low-dimensional subset of all edge tables. Kline's unrestricted edge model is the
whole edge table. Therefore the unrestricted edge model is more general.

The dimension count makes this precise. Suppose there are $J$ firms and $|E|$
observed directed edges. Holding the edge compositions $\mu_{jk}$ fixed, the
low-rank edge representation is governed by the shared parameters
\[
(\psi_1,\ldots,\psi_J),\qquad
(v_1,\ldots,v_J),\qquad
\Lambda,
\]
with at most
\[
J+Jd_v+d_ud_v
\]
raw degrees of freedom before normalizations. Kline's unrestricted edge vector
$\Delta$ has one free number per directed edge, so it has $|E|$ degrees of
freedom. When the mobility graph has many edges relative to the factor
dimensions, a generic vector in $\mathbb R^{|E|}$ cannot lie in the image of the
low-rank representation. Making the factors large enough to fit every edge would
turn the exercise into a saturated edge model, not the fixed low-rank wage
surface $(\star)$.

\subsubsection{Linear Conclusion}

The nesting answer is asymmetric:
\[
\text{low-rank wage surface}
\quad\Longrightarrow\quad
\text{structured Kline edge effects},
\]
but
\[
\text{unrestricted Kline edge effects}
\quad\not\Longrightarrow\quad
\text{low-rank wage surface}.
\]
Kline's $\Delta_{jk}$ is best understood as a diagnostic or statistical
interface. It records average wage changes on mover edges. Those edge changes
may be produced by AKM firm effects, low-rank interactions, match effects,
selection, or other mechanisms. But the unrestricted edge model does not identify
which mechanism generated them, and it does not require the shared-factor
restrictions imposed by the low-rank wage surface.

\section{Clean Takeaways}

\begin{enumerate}
\item The event study is mainly a parallel-trends / realized-mover diagnostic. It can
pass with time-invariant match effects because those match effects sit in wage levels
and treatment-effect heterogeneity, not necessarily in pre-trends.

\item Paired movers test separability more directly. For two workers making the same
$A\to B$ move, firm effects cancel, so the test asks whether wage gaps are portable
from origin to destination.

\item SW's additional low-power mechanism is pooling. Within-subsample event-study
coefficients can reveal problems, while between-sample mean differences push the
pooled coefficient back toward one.

\item Kline Section 3 says AKM is a path-independence restriction on unrestricted
directed edge effects. Cycle sums are the empirical content.

\item Kline's unrestricted edge model is not nested in the low-rank worker-firm wage
surface. AKM is nested as $\Lambda=0$, and a low-rank surface implies structured
edge effects, but arbitrary edge effects are a more general reduced form.
\end{enumerate}


\end{document}
